\section{Scope, parameter choices and limitations}\label{sec:discussion}

\subsection{Order of constants}

For the independent-noise proof, the constants can be chosen in the following order.
The buffer radius \(R=2000100\) is fixed. Choose a certificate
exponent \(A>100\), which suffices for the path count, and then
choose \(q>0\) small enough for
\eqref{eq:witness-prebound}. Choose \(B\) sufficiently large to
suppress the two clock-only terms in
\eqref{eq:bad-cost}, and finally \(\eta_0>0\) small enough for its
noise term. This fixes \(T=1000B\), the dependence divisor \(D_{\mathrm{dep}}\),
and all comparison constants. Only after these choices is \(L\)
required to be sufficiently large for the lifted windows and
epoch estimates. The resulting \(L_0\) can be astronomically large.
The proof does not give an operating threshold of practical size.
For Theorem~\ref{thm:general}, choose the coarse walk and density
constants, then the target superset cost, then \(B\), and then the
incident noise strength; the product mesh is sent to zero afterwards.
The resulting threshold is astronomically small and applies to
coherent amplitude \(|h|\), rather than \(h^2\).

\subsection{A wrapping syndrome that does not clear}

The clean-start qualification cannot be replaced by a deterministic
claim that every syndrome clears in a short time. To see this,
choose two different parallel coordinate lines winding in direction
1, at coordinates 0 and \(w\) in direction 2, with the other
coordinates fixed and \(1\le w<L\). Their mod-two sum is the
boundary of the periodic width-\(w\) ribbon of \((1,2)\)-plaquettes
between them: its interior edges cancel and the longitudinal ends
close by periodicity.

This is an attainable even-degree syndrome. At every occupied
vertex there is one incoming marked direction and one outgoing
marked direction. There is never a pair of incoming marked
directions, so every correction predicate fails. The syndrome
remains frozen under any noiseless attempt sequence.

An analogous example exists for the good-map comparison.
Initialize a complete coarse coordinate line with zeros and all
other bits with ones. At every site on the line the four coordinate
sets \(A_i\) contain the zero predecessor \(d=0\), while \(A_5\)
contains the zero predecessor \(d=e_1\) along the same periodic
line. Every minimum in \eqref{eq:comparison} is therefore zero,
so that line persists under every good update. Arbitrary bottom
zeros cannot be charged as fresh rare failed bins.

The overlap argument avoids this issue by using an older field
that itself arose from all-one data in the preceding epoch.
A disagreement across its overlap forces a long path in that
particular random clean-start field. Estimates are conditioned
only before the older window starts, never on an arbitrary
intermediate syndrome or comparison configuration.

\subsection{What is and is not implied}

Drechsler, Oreg and Aharonov's own physical model is independent
single-qubit Pauli noise, as specified in their Definitions 2.25,
2.37 and 7.7 \cite{DOA}. Theorem~\ref{thm:main} answers their
time-independent question in that model.
Theorem~\ref{thm:correlated} extends the conclusion to spacetime local
stochastic Pauli noise on unit time cells, including bounded-range
multi-qubit and adaptive faults with a law independent of the logical
input. Its description is \emph{Lindbladian correction,
trajectory-defined noise}; a general correlated law need not have
an additive Lindblad noise generator.

The local-stochastic condition depends on its time grid. A shifted
unit grid has parameter at most \(2\sqrt p\), and aligned \(k\)-fold
refinement has parameter at most \(p^{1/k}\), as proved in
Lemma~\ref{lem:grid}. These changes are not merely a fixed
rescaling of constants.

Their \emph{fully autonomous} notion asks for protection against
general local noise by a time-independent, local, translationally
invariant Lindbladian
\cite[Section 1.5.1, Table 1]{DOA}.
Theorem~\ref{thm:general} is the strongest result: exponential storage
under any sufficiently weak additive finite-range Markovian Lindblad
noise generator, with coherent and non-unital noise included.
It meets their fully autonomous definition for such noise, with
translation invariance when the noise is homogeneous.
The stochastic-Pauli theorems retain their stated structure and
correlated-history scope.
The thermal Davies-memory precedent is due to Alicki and the
Horodeckis \cite{AHHH2010}; the outside-equilibrium question is stated
explicitly in Drechsler--Oreg--Aharonov Section 1.4.1 \cite{DOA}.

The added noise generator must itself generate CPTP evolution.
Uniform correction-rate changes are handled by rescaling time;
small inhomogeneous rate changes can be treated above a common
positive baseline as additive local dissipators.
General perturbations of the correction's own jump operators,
including \(P\Pi\mapsto P\Pi+\delta K\), remain open, as does
coupling to a general non-Markovian quantum bath.
The former is part of the separate superoperator-stability problem
in their Section 1.8 \cite{DOA}.
No theorem here gives a general conversion of synchronous memories
to time-independent Lindbladians or addresses faulty final
syndrome measurement.

The final decoder may process the entire current syndrome. No
constant-time local classical implementation of that readout is
asserted. Its mathematical role is to define a fixed recovery
from the present configuration and to detect an invariant
homology class along small-component histories.

A logical membrane added to an initial error does not change
its syndrome, so syndrome geometry cannot determine that
initial logical class. From an encoded start it is known to
be trivial, and the history estimate prevents its subsequent
change with high probability. On \(\Z^4\), the result instead
states uniform rooted syndrome tails and vanishing low-noise
density. These are the exact finite- and infinite-volume
conclusions of the proof.
